% ============================================================================ INTRODUCTION
\section{Introduction}\label{sec:intro}

\IEEEPARstart{A}{rtificial-intelligence} and high-performance processors are approaching one
to two kilowatts per package at core voltages below one volt, so the package must deliver on
the order of a kiloampere~\cite{Krishnakumar2026DVPD,Ye2026TenKiloAmp,Swaminathan2026IVPDReview}.
At that current a single milliohm between the regulator and the load would dissipate about a
kilowatt, and moving voltage regulation into or directly beneath the package has become the
preferred scalable option~\cite{VPD2023Architectures,Ye2026TenKiloAmp}. Distributed vertical
power delivery (DVPD) goes one step further: many regulator cells, each beneath its own load
region, so that high voltage travels laterally and the kiloampere rail travels only
micrometers~\cite{Krishnakumar2026DVPD,Sriharini_ECTC_24,ISCAS2024DVPD}.

Bringing the regulator to the point of load changes the package it sits in. The regulator can be
embedded between the substrate and the die~\cite{Krishnakumar2026DVPD}, stacked on the die
backside~\cite{GoogleBacksideIVR2021}, attached to the package-substrate backside inside an
opening in the board~\cite{ArmCutoutVR2026}, or mounted on the board backside beneath the
package, as in current vertical-power modules~\cite{Vicor_VPD_AI,Baek2022LEGOPoL,
Radhakrishnan2021PowerDelivery,VPD_transactions_on_power_electronics_new}. Each choice decides
which rail crosses each vertical interconnect and at what current, where some 13\,\% of the input
power is dissipated inside the stack, which materials expand against each other, which solder
joints carry the deformation, and which vias must be manufactured. The converter-level decisions
are coupled to these as tightly: the regulator count fixes the footprint of each cell, the device
and passive selection fixes the component heights and ratings, and the assignment of components
to tiers fixes which connections must cross which dielectric depth. A design that is electrically
optimal can therefore exceed a via aspect-ratio rule, overload the ball field, outrun the cooling
solution, run its devices beyond their ratings or drive the critical solder joint harder than an
alternative.

Commercial multiphysics tools couple electrical, thermal and structural analysis of a given
layout, including stress, warpage and solder-life
estimates~\cite{Synopsys_RedHawkSC_ET,Cadence_Celsius,Ansys_Sherlock,Keysight_PIPro}; the
regulator, its tier order and its placement are inputs to them rather than design variables.
DVPD design frameworks have made the regulator a variable: the framework of~\cite{Krishnakumar2026DVPD}
co-optimizes switch sizing, regulator count, floorplan, interconnect allocation, IR drop and
transient droop for a given hybrid topology, and verifies the result with a self-consistent
electrothermal model with substrate-embedded microfluidic
cooling~\cite{Choi2025SelfConsistent,Choi2025AutoET}. Thermal comparisons of DVPD integration
variants~\cite{Choi2024ITherm,Choi2025EmbeddedCooling}, electrothermal co-optimization with
thermal TSVs~\cite{Wang2025TTSV} and regulator-placement studies for power integrity and
electromagnetic interference~\cite{Xu2019VRplacement,Hossen2020DSEarXiv,Zhang2019PDNbenchmark}
have also been reported. To our knowledge, no published study carries a converter synthesized
over its topology from discrete parts into a mounted, explicitly resolved package, compares its
integration architectures under identical cooling, and evaluates their thermomechanical and
solder-joint consequences together with the ratings of the selected parts and the
realizability of the vertical interconnect (Table~\ref{tab:prior}).

% Capability comparison with the closest prior work (evidence: review_work/novelty/notes.txt;
% the detailed version with quotes and pointers is Table S1 of the supplementary material).
\begin{table}[!t]
\caption{Capabilities stated in the closest prior work and in this work.}
\label{tab:prior}
\centering
\scriptsize
\setlength{\tabcolsep}{2.5pt}
\begin{tabular}{@{}lcccc@{}}
\toprule
 & \cite{Krishnakumar2026DVPD} & \cite{Choi2025SelfConsistent}
 & \cite{Synopsys_RedHawkSC_ET,Cadence_Celsius,Ansys_Sherlock} & This work \\
\midrule
Converter topology search           & --         & --\textsuperscript{a} & ?          & \checkmark\textsuperscript{b} \\
Discrete device/passive library & --         & --\textsuperscript{a} & ?          & \checkmark \\
Regulator-count optimization        & \checkmark & --                    & ?          & \checkmark\textsuperscript{c} \\
Floorplan, interconnect allocation  & \checkmark & --\textsuperscript{a} & ?          & $\circ$ \\
IR drop, transient droop            & \checkmark & ?                     & $\circ$    & -- \\
Package-explicit geometry           & $\circ$    & \checkmark            & \checkmark & \checkmark \\
Integration comparison & $\circ$    & --                    & ?          & \checkmark \\
Electrothermal coupling             & \checkmark & \checkmark            & \checkmark & \checkmark \\
Cooling boundary technology-neutral & -- & --                    & \checkmark & $\circ$\textsuperscript{d} \\
Thermomechanics                     & --         & ?                     & \checkmark & \checkmark \\
All-site joint search               & --         & ?                     & ?          & \checkmark \\
Local solder response               & --         & ?                     & \checkmark & \checkmark\textsuperscript{e} \\
Interconnect realizability          & $\circ$    & ?                     & ?          & \checkmark \\
Component-rating feasibility        & $\circ$    & ?                     & ?          & \checkmark \\
Per-stage reference verification    & $\circ$    & ?                     & $\circ$    & \checkmark\textsuperscript{f} \\
Runtime reported                    & --         & $\circ$               & ?          & \checkmark \\
\bottomrule
\end{tabular}
\par\smallskip
\parbox{\columnwidth}{\scriptsize
\checkmark: stated and demonstrated; $\circ$: partial; --: not part of the work; ?: not
determinable from the accessible text (full text of~\cite{Choi2025SelfConsistent} not
accessible; product pages of~\cite{Synopsys_RedHawkSC_ET,Cadence_Celsius,Ansys_Sherlock}
silent). \textsuperscript{a}\,Supplied by the electrical
framework~\cite{Krishnakumar2026DVPD}. \textsuperscript{b}\,All $(N,K)$ members within a
component budget; $K_A$ and $\mathbf{B}$ set by rule~\cite{Abdelzaher2025ECTC,
Abdelzaher2026EPEPS}. \textsuperscript{c}\,Electrical screen; the strongest candidates are
carried through the coupled chain. \textsuperscript{d}\,Declared film boundary;
demonstrated with microchannel plates. \textsuperscript{e}\,Inelastic response, no life.
\textsuperscript{f}\,Converter loss model validated on other
members~\cite{Abdelzaher2025ECTC,Abdelzaher2026EPEPS}; solder law by sensitivity.}
\end{table}


This paper presents such a framework and makes the following distinctions explicit. From the
DVPD framework of~\cite{Krishnakumar2026DVPD} it adopts the 1-kW, 48-to-1-V specification, the
embedded integration, the TSV and ball-grid technology, the routing model that scores the
regulator-count screen, and the input-current model used as one scenario. From our converter
work~\cite{Abdelzaher2025ECTC,Abdelzaher2026EPEPS} it adopts the generalized hybrid
switched-capacitor--buck family, its closed-form loss model and the selection of catalog
devices, passives and switching frequency. It extends the count optimization
of~\cite{Krishnakumar2026DVPD}, which sizes switch widths continuously for one topology, by
re-synthesizing the topology from discrete parts at every count, and it carries the
self-consistent electrothermal treatment of~\cite{Choi2025SelfConsistent}, applied there to one
embedded architecture, to an explicit mounted package with a declared cooling boundary. New are
the comparison of three placements of one synthesized regulator---embedded, package-backside and
board-backside---at a common temperature and under two identical coolings, after the candidates
that violate a component rating are filtered out; a thermoelastic solution of the mounted
assembly with a search over every ball site, a viscoplastic submodel of the critical joint and
an optional mechanical-to-thermal feedback through the top interface; and a per-net
realizability check of the tier order whose thermal and mechanical consequences are evaluated
on the same geometry. Every reduced model of the package chain is checked against a matched
reference, and one coupled evaluation takes \EvalHours{} on a shared workstation. Unlike
\cite{Krishnakumar2026DVPD}, the framework resolves neither spatial IR drop nor transient droop,
and the package result is not fed back into the synthesis automatically; the conclusions are
restricted to steady-state loss, temperature, deformation and joint response.

The framework is used in three ways, each demonstrated on a \SI{1}{kW}, \SIrange{48}{1}{V}
system: to \emph{evaluate} one design against its ratings and package rules, to \emph{compare}
integration architectures on conversion, routing, thermal and mechanical grounds, and to
\emph{explore} the design space---regulator count and converter topology, and the package
decisions on which every ranking depends. The coupled answers differ from the single-domain ones. At a common temperature the
package-backside regulator is marginally ahead of the embedded one. Under one top cold plate the
converter's thermal path reverses that order for ideal decoupling and ties the two once the
embedded regulator's decoupling is implemented, while the board-backside regulator exceeds the
ratings of its own devices. With an identical bottom plate the package-backside regulator leads
electrically and the embedded one mechanically, and the board-backside regulator falls behind both
once its 1-V redistribution is in the coupled solution. A flat cold plate on a warped package closes
a mechanical-to-thermal loop that leaves the package-backside regulator without a feasible operating
point, and the tier order chosen for the vias also sets the device temperature, its via advantage
conditional on a thin flying-capacitor technology.

The contributions are as follows.
\begin{enumerate}
\item A device-to-joint electrothermal--mechanical framework for package-integrated voltage
regulation, in which topology-level synthesis from a discrete library is propagated into an
explicit package and evaluated electrically, thermally, mechanically, against component ratings
and for interconnect realizability on one geometry, with each coupling direction classified and
the interface feedback available as an option (Sections~\ref{sec:framework}--\ref{sec:models}).
\item Physics and verification for that chain: a rail-aware conductor ledger, a voltage-fed
package network with the ball count as a solved variable, an electrothermal fixed point with a
physical cold-plate boundary and its hydraulic cost, and a thermoelastic package model driving a
local solder submodel with a mesh-independent response metric, each checked against a matched
reference, with a full-package cross-check and an uncertainty budget for every compared output
(Sections~\ref{sec:models}--\ref{sec:verification}).
\item A demonstration of the three modes on one specification: evaluation of a design with its
input decoupling implemented; comparison of three integration architectures at a common
temperature and under two identical coolings, with a feasibility filter and platform-inclusive
accounting; and exploration of regulator count and topology with coupled evaluation of the
strongest candidates, and of the package decisions that reverse each ranking
(Sections~\ref{sec:evaluate}--\ref{sec:explore}).
\item Design rules that follow from the coupled solution, stated with the conditions under
which they hold (Section~\ref{sec:discussion}).
\end{enumerate}
